\section{Dissipation of the unweighted gradient}
\label{rcut:sec:gradient}

The next step is a pointwise dissipation bound for the unweighted
gradient $U$, with a fixed contraction rate outside a set of exponentially
small Gibbs mass.  We first perform the deterministic calculation: completing
squares removes the second differences and leaves a quadratic form in the
first differences.  The error estimate retains the variance factors needed
when local fields are large.  We then apply the probability bounds proved in
Sections~\ref{rcut:sec:comparison} and~\ref{rcut:sec:scalar} to identify
the exceptional set.

Two row sums will control, respectively, the quadratic terms and the
cubic error.  For a symmetric matrix $J$ with zero diagonal, define
\[
 R_2(J)=\max_i\sum_j J_{ij}^2,
 \qquad
 R_3(J)=\max_i\sum_j |J_{ij}|^3.
\]
We begin at the all-positive configuration $\mathbf 1$.  In the quadratic
form below all local quantities are evaluated there:
$h_i=(J\mathbf 1)_i$, $m_i=\tanh h_i$,
$v_i=1-m_i^2$, and $w_i=1-m_i$.  Define
\begin{equation}
\begin{split}
 M_J(p)={}&\sum_{i\ne j}J_{ij}p_i v_jp_j
       +2\sum_{i\ne j}J_{ij}^2p_i m_jv_jp_j\\
       &+\sum_{i<j}
          \frac{J_{ij}^2(v_jp_i+v_ip_j)^2}{w_i+w_j},
       \qquad p\in\R^n.
\end{split}
\label{rcut:eq:gradient-form}
\end{equation}
For a finite matrix $J$, every denominator is positive.  The estimates
below remain uniform as any local field tends to infinity, including
when these denominators approach zero.

To transfer the calculation to an arbitrary configuration $x$, use
the sign change
\[
 J^x=\diag(x)J\diag(x),\qquad
 f^x(y)=f(x_1y_1,\ldots,x_ny_n),\qquad
 p_i^x=x_i d_i f(x).
\]
This sign change conjugates the generator for $J$ to the generator for
$J^x$.  It also gives $d_i f^x(\mathbf 1)=p_i^x$ and
$U(f^x)(\mathbf 1)=U(f)(x)$, while preserving the matrix controls:
$\|J^x\|=\|J\|$, $R_k(J^x)=R_k(J)$ for $k=2,3$.
The calculation at $\mathbf 1$ therefore suffices for every configuration.

\begin{lemma}[Pointwise gradient calculation]
\label{rcut:lem:gradient-algebra}
Let $f_t=S_tf_0$ for the heat-bath generator with a symmetric
zero-diagonal interaction matrix $J$.  At $\mathbf 1$, write
$p_i=d_if_t$, $r_{ij}=d_id_jf_t$, and $a_{ij}=d_im_j$.
Then
\begin{equation}
 \frac12(\partial_t-L)U(f_t)
 =-\|p\|^2
   +\sum_{i\ne j}p_i a_{ij}(p_j-2r_{ij})
   -\sum_{i\ne j}w_jr_{ij}^2.
\label{rcut:eq:gradient-identity}
\end{equation}
If $\max_{i,j}|J_{ij}|\le1/10$, then, at every configuration $x$,
\begin{equation}
 \frac12(\partial_t-L)U(f_t)(x)
 \le -U(f_t)(x)+M_{J^x}(p^x)
             +16R_3(J)U(f_t)(x).
\label{rcut:eq:gradient-local-bound}
\end{equation}
The constants do not depend on the local fields, on $f_0$, or on $t$.
\end{lemma}

\begin{proof}
We first derive the exact evolution identity, then estimate the field
differences before completing the squares.  The half-difference $d_if$
does not depend on spin $i$.  At the all-positive point, the product
rule reads
\[
 d_i(ab)=a\,d_ib+b(x^{i,-})\,d_ia.
\]
For $j\ne i$, the first difference after a flip is
$d_jf(x^{i,-})=p_j-2r_{ij}$.  The contribution with $j=i$ to
$d_iL f$ is $-p_i$: both $m_i$ and $p_i$ are independent of spin $i$.
Separating this term gives
\[
 d_iL f=L(d_if)-p_i
                   +\sum_{j\ne i}a_{ij}(p_j-2r_{ij}).
\]
Apply the square identity $L(g^2)-2gL g=2\Gamma(g)$ with $g=d_if$.
The equality $\Gamma(d_if)=\sum_{j\ne i}w_jr_{ij}^2$ then identifies
the second-difference energy, proving
Equation~\eqref{rcut:eq:gradient-identity}.

To use the identity uniformly in the local fields, we must keep the
conditional variance in the estimate for $a_{ij}$.  Set $z=J_{ij}$
and $m=m_j(\mathbf 1)$.  Flipping spin $i$ changes the $j$-field by
$-2z$.  The hyperbolic-tangent subtraction formula therefore gives
\begin{equation}
 a_{ij}=v_j q(z,m),\qquad
 q(z,m)=\frac{\tanh(2z)}{2(1-m\tanh(2z))}.
\label{rcut:eq:gradient-field-difference}
\end{equation}
For real arguments with $|z|\le1/10$ and $|m|\le1$, the factor
$1-m\tanh(2z)$ is bounded below by $4/5$.
The bound is uniform as $m$ approaches either endpoint of $[-1,1]$,
so the variance factor can be retained without a field cutoff.

To track the constants in the cubic remainder, write $T=\tanh(2z)$
and expand the denominator exactly:
\[
 q(z,m)=\frac T2+\frac{mT^2}{2}
                   +\frac{m^2T^3}{2(1-mT)}.
\]
The inequalities $|T|\le2|z|$ and
$|\tanh u-u|\le |u|^3/3$ imply
\begin{equation}
 \begin{split}
 |q(z,m)-z-2mz^2|&\le7|z|^3,\\
 |q(z,m)-z|&\le\frac{27}{10}z^2,
 \qquad |q(z,m)|\le\frac54|z|.
 \end{split}
\label{rcut:eq:gradient-q-bounds}
\end{equation}
For the first estimate, the three terms of this decomposition contribute
errors bounded by $(4/3)|z|^3$, $(16/3)|z|^4$, and $5|z|^3$,
respectively.  Their sum is less than $7|z|^3$.
The second estimate is a consequence of the first, while the last is
immediate from Equation~\eqref{rcut:eq:gradient-field-difference}.
The hyperbolic-tangent inequality follows by integrating
$1-(\tanh u)'=\tanh^2u\le u^2$.
Returning to the field difference gives
\[
 a_{ij}=J_{ij}v_j+2J_{ij}^2m_jv_j+e_{ij},
 \qquad |e_{ij}|\le7|J_{ij}|^3v_j,
\]
with the variance factor retained in the remainder.

We can now eliminate the second differences.  Since $r_{ij}=r_{ji}$,
the terms for one unordered pair in
Equation~\eqref{rcut:eq:gradient-identity} combine as
\[
 -2(a_{ij}p_i+a_{ji}p_j)r_{ij}-(w_i+w_j)r_{ij}^2.
\]
Completing the square bounds this expression by
$(a_{ij}p_i+a_{ji}p_j)^2/(w_i+w_j)$.
We compare its numerator with the leading interaction term.  For the
same pair, write
\[
 A=a_{ij}p_i+a_{ji}p_j,\quad
 B=J_{ij}(v_jp_i+v_ip_j),\quad
 S=v_j|p_i|+v_i|p_j|.
\]
Equation~\eqref{rcut:eq:gradient-q-bounds} gives
\[
 |A-B|\le\frac{27}{10}|J_{ij}|^2S,
 \qquad |A|+|B|\le\frac94|J_{ij}|S.
\]
Hence $|A^2-B^2|\le7|J_{ij}|^3S^2$.
The remaining issue is the denominator: it can be small at a nearly
saturated field.  Weighted Cauchy--Schwarz keeps the variance factors
attached to the corresponding weights and gives
\begin{equation}
 \frac{S^2}{w_i+w_j}
 \le\frac{v_j^2}{w_j}p_i^2+\frac{v_i^2}{w_i}p_j^2,
 \qquad
 \frac{v_i^2}{w_i}=(1-m_i)(1+m_i)^2\le\frac{32}{27}.
\label{rcut:eq:gradient-denominator-control}
\end{equation}
This estimate also describes the limiting fields.  As $m_j\uparrow1$,
the contribution is at most $4(1-m_j)p_i^2$; as $m_j\downarrow-1$,
it is at most $2(1+m_j)^2p_i^2$.  Both bounds hold without comparing
the coupling size with the distance from saturation.

Only the accumulated replacement errors remain.  Replacing
$\sum_{i\ne j}p_i a_{ij}p_j$ by the first two terms of $M_J(p)$
costs at most
\[
 7\sum_{i\ne j}|J_{ij}|^3|p_i||p_j|
 \le7R_3(J)\|p\|^2.
\]
The last inequality uses symmetry and $2|p_i||p_j|\le p_i^2+p_j^2$.
For the completed squares,
Equation~\eqref{rcut:eq:gradient-denominator-control} bounds the total
replacement error by
\[
 \frac{224}{27}\sum_{i<j}|J_{ij}|^3(p_i^2+p_j^2)
 \le\frac{224}{27}R_3(J)\|p\|^2.
\]
Adding the two errors and using $7+224/27<16$ proves
Equation~\eqref{rcut:eq:gradient-local-bound} at $\mathbf 1$.
Sign conjugation then gives the asserted inequality at every $x$.
\end{proof}

The local calculation reduces dissipation to an upper bound on the
quadratic form.  We first give a bound valid at every configuration; it
will also control possible growth inside the exceptional set.
For every symmetric zero-diagonal $J$ and every $p$,
\begin{equation}
 M_J(p)\le
 \left[\|J\|+
       \left(\frac{4}{3\sqrt3}+\frac{32}{27}\right)R_2(J)
 \right]\|p\|^2.
\label{rcut:eq:rough-gradient}
\end{equation}
To see this, treat the three terms in
Equation~\eqref{rcut:eq:gradient-form} separately.  The first is bounded
by $\|J\|\|p\|^2$ because $0\le v_i\le1$.
For the second, the symmetric matrix $J\circ J$, with entries
$J_{ij}^2$, has norm at most $R_2(J)$; combine this with
$\max_{|m|\le1}|m(1-m^2)|=2/(3\sqrt3)$.
For the final term, the weighted Cauchy--Schwarz estimate in
Equation~\eqref{rcut:eq:gradient-denominator-control} gives
$(32/27)R_2(J)\|p\|^2$.  This proves
Equation~\eqref{rcut:eq:rough-gradient} uniformly in the local fields.

\begin{proposition}[Typical-state gradient dissipation]
\label{rcut:prop:typical-gradient}
Fix $0\le\beta<1/2$.  There are constants $c>0$ and $C<\infty$,
and events of disorder probability tending to one, on which a set
$\mathcal B_J\subseteq\{-1,1\}^n$ satisfies
\[
 \mu_J(\mathcal B_J)\le e^{-cn}
\]
and, for every $f_0$, every $t\ge0$, and every configuration $x$,
\begin{equation}
 (\partial_t-L)U(S_tf_0)(x)
 \le[-c+C\1_{\mathcal B_J}(x)]\,U(S_tf_0)(x).
\label{rcut:eq:gradient-potential}
\end{equation}
The constants depend only on the fixed $\beta$.  If
$0\le\beta\le7/20$, the exceptional set can be empty.
\end{proposition}

\begin{proof}
We now combine the deterministic calculation with the disorder estimates.
For the smaller temperatures, the rough bound already contracts at every
configuration.  For the remaining temperatures, the planted comparison
will provide contraction outside a small Gibbs set.  Begin on the
events of Proposition~\ref{rcut:prop:good-disorder}, where the deterministic
tolerances give
\[
 \begin{gathered}
 \|J\|\le2\beta+o(1),\qquad R_2(J)\le\beta^2+o(1),\\
 R_3(J)\le\max_{i,j}|J_{ij}|\,R_2(J)=o(1).
 \end{gathered}
\]
Sign conjugation preserves each bound, so they hold for every $J^x$.
The remainder in Lemma~\ref{rcut:lem:gradient-algebra} consequently tends
to zero uniformly over configurations, test functions, and times.

If $\beta\le7/20$, it suffices to substitute the disorder estimates
into Equation~\eqref{rcut:eq:rough-gradient}.  Its coefficient is at most
\[
 2\beta+
 \left(\frac{4}{3\sqrt3}+\frac{32}{27}\right)\beta^2+o(1)
 <\frac{189}{200}+o(1)<1.
\]
The rational bound follows from $4/(3\sqrt3)+32/27<2$ at
$\beta=7/20$, since both terms increase with nonnegative $\beta$.
There is therefore a fixed margin below one.  Absorbing the vanishing
remainder in Lemma~\ref{rcut:lem:gradient-algebra} proves
Equation~\eqref{rcut:eq:gradient-potential} with a fixed positive $c$
and $\mathcal B_J=\varnothing$ for all sufficiently large $n$.
The endpoint $\beta=0$ can also be read directly from
Equation~\eqref{rcut:eq:gradient-identity}: $a_{ij}=0$, $w_i=1$, and
$(\partial_t-L)U\le-2U$.  This part of the proof therefore uses no
nondegenerate Gaussian field law.

For $7/20<\beta<1/2$, we need a sharper estimate on most
configurations.  The inputs are proved in
Sections~\ref{rcut:sec:comparison} and~\ref{rcut:sec:scalar}, independently
of Proposition~\ref{rcut:prop:typical-gradient}.
First, Proposition~\ref{rcut:prop:scalar-gap} bounds the normalized
variational expression in Equation~\eqref{rcut:eq:variational} by $0.9969$.
Choose $\delta=0.0031$ and $\rho=\delta/2$ to retain half this strict
margin.  Proposition~\ref{rcut:prop:planted-comparison} then gives
$\kappa>0$ for which the event
\[
 \mathcal A_n=
 \left\{J:\sup_{\|p\|=1}M_J(p)>1-\rho\right\}
\]
has planted probability at most $e^{-\kappa n}$ for all sufficiently
large $n$.  Under the planted law, the entries above the diagonal are
independent with mean $\beta^2/n$ and variance $\beta^2/n$;
Section~\ref{rcut:sec:comparison} defines and analyzes this law.
To translate the matrix event into a set of spin configurations, define
\begin{equation}
 \mathcal B_J=\{x:J^x\in\mathcal A_n\}.
\label{rcut:eq:gradient-bad-set}
\end{equation}
Lemma~\ref{rcut:lem:annealed-transfer} now yields
$\mu_J(\mathcal B_J)\le e^{-\kappa n/2}$ on events whose original
disorder probability tends to one.  Intersect these with the good
disorder events used above.  The supremum over all unit vectors in the
definition ensures that the set in
Equation~\eqref{rcut:eq:gradient-bad-set} depends only on $J$, rather
than on the evolving function.

The Gibbs mass estimate is established.  To obtain the pointwise
dissipation bound, take $n$ large enough that $16R_3(J)\le\rho/2$.
For $x\notin\mathcal B_J$, homogeneity of $M_{J^x}$ and
Lemma~\ref{rcut:lem:gradient-algebra} leave the fixed margin
\[
 \frac12(\partial_t-L)U(S_tf_0)(x)
 \le-\frac\rho2 U(S_tf_0)(x).
\]
At every configuration, including those in the exceptional set,
Equation~\eqref{rcut:eq:rough-gradient} and
Lemma~\ref{rcut:lem:gradient-algebra} give $(\partial_t-L)U\le K U$
for some nonnegative finite $K$ depending only on $\beta$.
For instance, the eventual estimates $\|J\|\le2$, $R_2(J)\le1$,
and $16R_3(J)\le1$ permit $K=8$.
Choose $c=\min\{\rho,\kappa/2\}$ and $C=K+c$ to combine the Gibbs
mass and dissipation bounds on the same event.  Since the local
calculation applies to every function, this event and these constants
work simultaneously for all $f_0$, $t$, and $x$.
\end{proof}
